\apartado{}{Mezcla}{preambulo}
\bajada{Acá no se dice de qué apartado es cada ejercicio. Ésa es toda la
dificultad, y es la que van a enfrentar en un parcial y en el trabajo: nadie
les va a avisar qué herramienta usar.}

\minihistoria{Por qué esta sección existe}{%
  Resolver ocho ejercicios seguidos de combinatoria no prueba que entendieron
  combinatoria: prueba que reconocieron el molde de la página anterior.\par
  Si acá elegís bien la herramienta, aprendiste. Si te trabás, no significa que
  no sepas — significa que hasta ahora el enunciado te venía diciendo la mitad
  de la respuesta.}

\ejercicio{4}{En una investigación se aplicaron 5 pruebas a cada participante y
  el equipo quiere contrabalancear el orden. Tienen 40 participantes.
  \begin{enumerate}[label=(\alph*),topsep=2pt,itemsep=1pt]
    \item ¿Cuántos órdenes posibles existen?
    \item ¿Alcanzan los 40 participantes para un contrabalanceo completo?
    \item Si no alcanzan, ¿qué proponés?
  \end{enumerate}}{}
\renglones[6]
\respuesta{(a) $5!=120$. (b) No: harían falta 120. (c) Respuesta abierta
(cuadrado latino de 5 secuencias, 8 participantes por secuencia).}
\solucion{%
  (a) $5! = 120$ órdenes.
  \par
  (b) \textbf{No alcanzan.} Para el contrabalanceo completo harían falta 120
  participantes, o un múltiplo.
  \par
  (c) Lo estándar es un \textbf{cuadrado latino}: elegir 5 secuencias en las
  que cada prueba aparezca exactamente una vez en cada posición, y asignar 8
  participantes a cada una. Se pierde información sobre interacciones entre
  órdenes, pero se controla el efecto de posición, que es lo que importa.}

\ejercicio{3}{Un servicio atendió a 180 personas el año pasado. De ellas, 54
  presentaron sintomatología ansiosa y 36 sintomatología depresiva; 18
  presentaron las dos. ¿Qué probabilidad hay de que una persona elegida al azar
  de ese grupo presente al menos una de las dos?}{}
\renglones[4]
\respuesta{$72/180 = 0{,}40$.}
\solucion{%
  \[ P(A\cup D)=P(A)+P(D)-P(A\cap D)
     =\frac{54}{180}+\frac{36}{180}-\frac{18}{180}=\frac{72}{180}=0{,}40 \]
  \par
  En personas: $54+36-18=72$. Sin restar los 18 que están en los dos grupos
  darías 90, contándolos dos veces — un 25\% de más.
  \par
  \textit{Es del apartado 2.5, que todavía no vimos en este cuaderno. Si te
  salió igual razonando con el diagrama, es buena señal: la regla de la suma
  es un dibujo antes que una fórmula.}}

\ejercicio{2}{Dos psicólogos codifican 80 registros de observación y coinciden
  en 60. Cada uno marca «conducta presente» en el 50\% de los casos.
  \begin{enumerate}[label=(\alph*),topsep=2pt,itemsep=1pt]
    \item ¿Cuál es el acuerdo observado?
    \item ¿Cuánto se esperaría por azar?
    \item ¿Es mejor o peor acuerdo que el del ejercicio de las 50 entrevistas
      ($P_o=0{,}80$, $P_e=0{,}58$)?
  \end{enumerate}}{}
\renglones[6]
\respuesta{(a) $75\%$. (b) $50\%$. (c) Mejor: $\kappa=0{,}50$ contra $0{,}52$…
en realidad casi igual.}
\solucion{%
  (a) $60/80 = 0{,}75 = 75\%$.
  \par
  (b) Cada uno dice «sí» la mitad de las veces:
  $0{,}5\times 0{,}5 + 0{,}5\times 0{,}5 = 0{,}50$.
  \par
  (c) $\kappa = \dfrac{0{,}75-0{,}50}{1-0{,}50} = \dfrac{0{,}25}{0{,}50}
  = 0{,}50$, contra $0{,}52$ del otro caso.
  \par
  \textbf{Lo interesante}: el acuerdo crudo es menor ($75\%$ contra $80\%$) y
  sin embargo el kappa es prácticamente el mismo. Porque acá la categoría es
  frecuente —la mitad de los casos— y el azar regala mucho menos. Comparar
  porcentajes de acuerdo entre estudios con categorías de distinta frecuencia
  no significa nada.}

\ejercicio{2}{Un test de screening de consumo problemático tiene 6 ítems, cada
  uno con 3 opciones. ¿Cuántas formas distintas hay de responderlo? ¿Y cuántos
  puntajes distintos puede dar, si cada opción vale 0, 1 o 2?}{}
\renglones[4]
\respuesta{$3^{6}=729$ formas; 13 puntajes (de 0 a 12).}
\solucion{%
  Formas de responder: $3^{6} = 729$.
  \par
  Puntajes: de $6\times 0 = 0$ a $6\times 2 = 12$, o sea \textbf{13} puntajes
  posibles.
  \par
  Es el mismo par de preguntas del apartado 2.1 con otros números, y sirve para
  lo mismo: 729 formas distintas se comprimen en 13 puntajes, así que muchísimas
  personas con perfiles distintos van a compartir el mismo total.}

\ejercicio{3}{En una muestra de 300 adolescentes, un tamizaje de riesgo
  suicida dio positivo en 24 casos. De esos 24, la entrevista confirmó riesgo
  en 9. Además, entre los 276 que dieron negativo, la entrevista detectó riesgo
  en 3 que se habían escapado.
  \begin{enumerate}[label=(\alph*),topsep=2pt,itemsep=1pt]
    \item Armá la tabla de $2\times 2$.
    \item Calculá sensibilidad, especificidad y VPP.
    \item ¿Cuál de los tres le informarías a un adolescente que dio positivo?
  \end{enumerate}}{}
\renglones[8]
\respuesta{$VP=9$, $FP=15$, $FN=3$, $VN=273$. Sens.\ $=9/12=0{,}75$;
espec.\ $=273/288\approx 0{,}948$; VPP $=9/24=0{,}375$. Se informa el VPP.}
\solucion{%
  (a) Con riesgo real: $9+3=12$. Sin riesgo: $300-12=288$.
  \begin{center}\small
  \begin{tabular}{@{}lccc@{}}
  \toprule
   & \textbf{Riesgo sí} & \textbf{Riesgo no} & \textbf{Total}\\
  \midrule
  \textbf{Test $+$} & 9 & 15 & 24\\
  \textbf{Test $-$} & 3 & 273 & 276\\
  \textbf{Total} & 12 & 288 & 300\\
  \bottomrule
  \end{tabular}
  \end{center}
  \par
  (b) Sensibilidad $=\dfrac{9}{12}=0{,}75$; \;
  especificidad $=\dfrac{273}{288}\approx 0{,}948$; \;
  VPP $=\dfrac{9}{24}=0{,}375$.
  \par
  (c) El \textbf{VPP}: es la única que responde «di positivo, ¿qué significa
  para mí?». Y da $37{,}5\%$ — menos de la mitad.
  \par
  \textbf{Y sin embargo el instrumento sirve.} Con riesgo suicida, un falso
  negativo es catastrófico y una falsa alarma cuesta una conversación. Un VPP
  bajo es el precio deliberado de una sensibilidad alta, y en este caso es un
  precio que vale la pena pagar. Que un número sea bajo no lo vuelve un
  problema: hay que preguntarse qué se compró con él.}
\enclase{Este ejercicio junta 2.3 completo y obliga a construir la tabla desde
un texto, que es como llegan los datos en la realidad. Si sólo hacen uno de la
mezcla, que sea éste.}

\ejercicio{1}{De 250 pacientes atendidos, 75 no completaron el tratamiento.
  ¿Qué probabilidad hay de que un paciente elegido al azar lo haya completado?
  ¿Qué tipo de probabilidad es?}{}
\renglones[3]
\respuesta{$175/250 = 0{,}70$. Frecuentista.}
\solucion{%
  $P(\text{completó}) = \dfrac{250-75}{250} = \dfrac{175}{250} = 0{,}70$, o por
  complemento: $1 - \dfrac{75}{250} = 1-0{,}30 = 0{,}70$.
  \par
  Es \textbf{frecuentista}: sale de contar lo observado.}

\ejercicio{4}{Un equipo de investigación tiene 12 participantes y quiere
  formar 3 grupos de 4 para una intervención grupal. Un integrante propone
  calcular $C(12,4) = 495$ y dice que ésas son las formas posibles de armar los
  grupos.
  \textbf{¿Está bien?} Explicá qué cuenta realmente el 495, y si alcanza para
  responder la pregunta que se hicieron.}{}
\renglones[6]
\respuesta{No alcanza: 495 cuenta sólo las formas de armar \emph{el primer}
grupo de 4, no los tres grupos completos.}
\solucion{%
  El $C(12,4)=495$ está bien calculado, pero cuenta \textbf{una sola cosa}: de
  cuántas maneras se elige el primer grupo de 4 entre los 12.
  \par
  Para los tres grupos hay que seguir: elegido el primero quedan 8, y
  $C(8,4)=70$; los 4 restantes forman el tercero sin elección. Eso da
  $495\times 70 = 34\,650$.
  \par
  Y todavía falta un detalle: si los tres grupos son \emph{intercambiables}
  —ninguno tiene un rol distinto— cada repartición quedó contada $3!=6$ veces,
  así que serían $34\,650/6 = 5\,775$.
  \par
  \textbf{Lo importante no es el número}, es que la primera respuesta parecía
  completa y respondía menos de un tercio de la pregunta. Cuando un problema
  tiene varias etapas, un solo $C(n,k)$ casi nunca alcanza.}
\enclase{Es el ejercicio más difícil del cuaderno. No esperes que lo resuelvan
entero: alcanza con que detecten que 495 no puede ser la respuesta completa.}
